\documentclass[11pt]{article}

\usepackage[T1]{fontenc}
\usepackage{lmodern}
\usepackage{amsmath,amssymb,amsthm,mathtools}
\usepackage[margin=1.08in]{geometry}
\usepackage{microtype}
\usepackage[colorlinks=true,linkcolor=blue,citecolor=blue,urlcolor=blue]{hyperref}

\newtheorem{theorem}{Theorem}[section]
\newtheorem{proposition}[theorem]{Proposition}
\newtheorem{corollary}[theorem]{Corollary}
\theoremstyle{definition}
\newtheorem{definition}[theorem]{Definition}
\newtheorem{problem}[theorem]{Problem}
\newtheorem{conjecture}[theorem]{Conjectural sharp value}
\theoremstyle{remark}
\newtheorem{remark}[theorem]{Remark}

\newcommand{\C}{\mathbb C}
\newcommand{\N}{\mathbb N}
\newcommand{\BH}{\mathcal B(\mathcal H)}
\newcommand{\AR}{\mathbb A_R}
\newcommand{\QAR}{\mathbb Q\mathbb A_R}
\newcommand{\Rat}{\operatorname{Rat}}
\newcommand{\cb}{\mathrm{cb}}
\newcommand{\fd}{\mathrm{fd}}
\newcommand{\Mat}{\mathrm M}
\newcommand{\id}{I}
\newcommand{\spec}{\sigma}

\title{The Quantum Annulus Spectral Constant Problem}
\author{}
\date{Status as of July 23, 2026}

\begin{document}

\maketitle

\begin{abstract}
Fix \(R>1\).  The quantum annulus of type \(R\) is the class of
invertible Hilbert-space operators \(T\) satisfying
\(\|T\|\leq R\) and \(\|T^{-1}\|\leq R\).  The associated spectral
constant \(K(R)\) is the least universal constant for which
\[
 \|f(T)\|
 \leq K(R)\max_{R^{-1}\leq |z|\leq R}|f(z)|
\]
holds for every operator in this class and every rational function
with no pole on the closed annulus.  This note gives a self-contained
formulation of the problem, its scaling normalization, its scalar,
finite-dimensional, and completely bounded variants, the known
bounds, and exact criteria for a solution.  The sharp value for a
fixed finite \(R>1\) remains unknown.  The natural candidate is
\(K(R)=2\).
\end{abstract}

\section{The annulus and its quantum operator class}

For \(R>1\), write
\[
 \AR=\{z\in\C:R^{-1}<|z|<R\},
 \qquad
 \overline{\AR}=\{z\in\C:R^{-1}\leq |z|\leq R\}.
\]
If \(\mathcal H\) is a complex Hilbert space, define
\begin{equation}
 \QAR(\mathcal H)
 =
 \left\{
 T\in\mathcal B(\mathcal H):
 T\text{ is invertible},\
 \|T\|\leq R,\
 \|T^{-1}\|\leq R
 \right\}.
 \label{eq:quantum-annulus}
\end{equation}
When the underlying Hilbert space is understood, it will be suppressed
from the notation.

Every \(T\in\QAR(\mathcal H)\) satisfies
\[
 \spec(T)\subseteq\overline{\AR}.
\]
Indeed, \(|\lambda|\leq R\) for \(\lambda\in\spec(T)\), while
\(\lambda^{-1}\in\spec(T^{-1})\) gives \(|\lambda|\geq R^{-1}\).

The following positive-operator condition is another useful
description of the class \cite{McCulloughPascoe2023}.

\begin{proposition}[Defect characterization]
\label{prop:defect-characterization}
For an invertible \(T\in\mathcal B(\mathcal H)\), the following are
equivalent:
\begin{enumerate}
 \item \(T\in\QAR(\mathcal H)\);
 \item
 \begin{equation}
  \beta_R(T^*,T)
  :=
  (R^2+R^{-2})\id
  -T^*T-(T^*T)^{-1}
  \succeq0.
  \label{eq:beta-defect}
 \end{equation}
\end{enumerate}
\end{proposition}

An invertible operator satisfying \(\beta_R(T^*,T)=0\) is called a
\emph{quantum annulus unitary}.  Such an operator need not be normal;
this terminology should not be confused with a normal operator whose
spectrum lies on the boundary of the classical annulus.

\subsection{Reduction of a general annulus to the symmetric form}

Let \(0<a<b\), and set
\[
 \mathbb A(a,b)=\{z\in\C:a<|z|<b\}.
\]
The natural quantum class associated with this annulus is
\[
 \mathbb Q\mathbb A(a,b;\mathcal H)
 =
 \left\{
 T\in\mathcal B(\mathcal H):
 T\text{ is invertible},\
 \|T\|\leq b,\
 \|T^{-1}\|\leq a^{-1}
 \right\}.
\]
Put
\begin{equation}
 c=\sqrt{ab},
 \qquad
 R=\sqrt{\frac ba},
 \qquad
 S=c^{-1}T.
 \label{eq:scaling-parameters}
\end{equation}
Then
\[
 T\in\mathbb Q\mathbb A(a,b;\mathcal H)
 \quad\Longleftrightarrow\quad
 S\in\QAR(\mathcal H),
\]
and \(z\mapsto c^{-1}z\) maps \(\mathbb A(a,b)\) onto \(\AR\).
Moreover, if \(g(w)=f(cw)\), then \(f(T)=g(S)\) and
\[
 \max_{a\leq|z|\leq b}|f(z)|
 =
 \max_{R^{-1}\leq|w|\leq R}|g(w)|.
\]
Consequently, the spectral constant for a general concentric annulus
depends only on the ratio \(b/a\), equivalently on its conformal
modulus.  It is therefore enough to study the symmetric one-parameter
family \(\AR\).

\begin{remark}[A second convention]
Some papers use \(0<\rho<1\) and write the symmetric annulus as
\(\{\rho<|z|<\rho^{-1}\}\), with
\(\|\rho T\|,\|\rho T^{-1}\|\leq1\).  This is exactly the convention
above with \(\rho=R^{-1}\).
\end{remark}

\section{Spectral and complete spectral constants}

For a compact set \(X\subseteq\C\), let \(\Rat(X)\) be the algebra of
rational functions whose poles lie outside \(X\), equipped with
\[
 \|f\|_X=\max_{z\in X}|f(z)|.
\]
If \(\spec(T)\subseteq X\), then \(f(T)\) is defined by the rational
functional calculus.

\begin{definition}[\(K\)-spectral set]
\label{def:spectral-set}
A compact set \(X\subseteq\C\) is a \(K\)-spectral set for \(T\) if
\[
 \|f(T)\|\leq K\|f\|_X
 \qquad
 \bigl(f\in\Rat(X)\bigr).
\]
For a fixed \(T\in\QAR(\mathcal H)\), define
\[
 \kappa_R(T)
 =
 \sup_{\substack{f\in\Rat(\overline{\AR})\\
                  \|f\|_{\overline{\AR}}\leq1}}
 \|f(T)\|.
\]
\end{definition}

\begin{definition}[Universal scalar constant]
\label{def:universal-constant}
The spectral constant of the quantum annulus is
\begin{equation}
 K(R)
 =
 \sup_{\mathcal H}
 \ \sup_{T\in\QAR(\mathcal H)}
 \kappa_R(T).
 \label{eq:K-def}
\end{equation}
Equivalently, \(K(R)\) is the least constant \(K\) such that
\(\overline{\AR}\) is a \(K\)-spectral set for every
\(T\in\QAR(\mathcal H)\), on every complex Hilbert space
\(\mathcal H\).
\end{definition}

There is also a matrix-valued version.  If
\(F=[f_{ij}]\in\Mat_m(\Rat(X))\), define
\[
 F(T)=[f_{ij}(T)]
 \quad\text{on }\mathcal H^{\oplus m},
 \qquad
 \|F\|_X=\max_{z\in X}\|F(z)\|_{\Mat_m}.
\]

\begin{definition}[Universal complete constant]
\label{def:complete-constant}
The completely bounded, or complete spectral, constant is
\begin{equation}
 K_{\cb}(R)
 =
 \sup_{\mathcal H}
 \ \sup_{T\in\QAR(\mathcal H)}
 \ \sup_{m\geq1}
 \ \sup_{\substack{F\in\Mat_m(\Rat(\overline{\AR}))\\
                   \|F\|_{\overline{\AR}}\leq1}}
 \|F(T)\|.
 \label{eq:Kcb-def}
\end{equation}
Necessarily,
\[
 K(R)\leq K_{\cb}(R).
\]
\end{definition}

\subsection{Finite-dimensional constants}

For \(n\geq1\), set
\[
 K_n(R)
 =
 \sup_{\substack{T\in\Mat_n(\C)\ {\rm invertible}\\
                  \|T\|,\|T^{-1}\|\leq R}}
 \kappa_R(T),
 \qquad
 K_{\fd}(R)=\sup_{n\geq1}K_n(R).
\]
Then
\begin{equation}
 1\leq K_n(R)\leq K_{\fd}(R)\leq K(R)\leq K_{\cb}(R).
 \label{eq:constant-hierarchy}
\end{equation}
The distinction between \(K_{\fd}(R)\) and \(K(R)\) is substantive:
some of the strongest matrix estimates currently use
finite-dimensional extremal-function arguments whose validity on an
arbitrary Hilbert space is not known.

\section{The sharp-constant problem}

\begin{problem}[General quantum annulus spectral constant]
\label{prob:general-constant}
For each fixed \(R>1\), determine the exact functions
\[
 R\longmapsto K(R)
 \qquad\text{and}\qquad
 R\longmapsto K_{\cb}(R).
\]
In particular:
\begin{enumerate}
 \item determine whether \(K(R)=2\) for every \(R>1\);
 \item determine whether \(K_{\cb}(R)=2\) for every \(R>1\);
 \item determine whether
 \[
  K_{\fd}(R)=K(R)=K_{\cb}(R);
 \]
 \item identify extremal or asymptotically extremal pairs
 \((T,f)\), and decide whether extremality can always be witnessed in
 finite dimension.
\end{enumerate}
\end{problem}

\begin{conjecture}
\label{conj:sharp-two}
For every \(R>1\),
\begin{equation}
 K(R)=2.
 \label{eq:conjectural-two}
\end{equation}
The stronger complete form asks whether \(K_{\cb}(R)=2\).
\end{conjecture}

Equation \eqref{eq:conjectural-two} is the natural sharp-value
candidate because \(2\) is already a universal lower bound and is the
exact limiting value as the annulus becomes thick.  It is not
currently a theorem for a general fixed finite \(R>1\).

\subsection{Equivalent Laurent-polynomial targets}

Laurent polynomials are uniformly dense in
\(\Rat(\overline{\AR})\).  Therefore the scalar assertion
\(K(R)\leq2\) is equivalent to the following concrete family of
inequalities:
\begin{problem}[Sharp scalar Laurent inequality]
\label{prob:scalar-laurent}
Prove that, for every complex Hilbert space \(\mathcal H\), every
\(T\in\QAR(\mathcal H)\), every \(N\geq0\), and every
\(a_{-N},\ldots,a_N\in\C\),
\begin{equation}
 \left\|
  \sum_{k=-N}^{N}a_kT^k
 \right\|
 \leq
 2
 \max_{R^{-1}\leq|z|\leq R}
 \left|
  \sum_{k=-N}^{N}a_kz^k
 \right|.
 \label{eq:scalar-laurent-target}
\end{equation}
\end{problem}

Likewise, the complete assertion \(K_{\cb}(R)\leq2\) is equivalent to
the matrix-valued version:
\begin{problem}[Sharp matricial Laurent inequality]
\label{prob:matrix-laurent}
Prove that, for every \(m,N\geq1\), every
\(A_{-N},\ldots,A_N\in\Mat_m(\C)\), and every
\(T\in\QAR(\mathcal H)\),
\begin{equation}
 \left\|
  \sum_{k=-N}^{N}A_k\otimes T^k
 \right\|
 \leq
 2
 \max_{R^{-1}\leq|z|\leq R}
 \left\|
  \sum_{k=-N}^{N}A_kz^k
 \right\|_{\Mat_m}.
 \label{eq:matrix-laurent-target}
\end{equation}
\end{problem}

The negative powers in
\eqref{eq:scalar-laurent-target}--\eqref{eq:matrix-laurent-target}
are essential.  A polynomial estimate alone does not solve the
annulus problem, because the rational functional calculus must see
both boundary components.

\section{What is known}

Define
\begin{equation}
 K_{\mathrm{Tom}}(R)
 =
 2\left(
  1+
  \frac{2R^2}
  {(R^2+1)\sqrt{R^4-1}}
 \right).
 \label{eq:Tomar-bound}
\end{equation}
Also set
\begin{equation}
 K_{\mathrm{BBC}}(R)
 =
 2+
 \min\left\{
 \frac{R+1}{\sqrt{R^2+R+1}},\,
 \sqrt{\frac{R^2+1}{R^2-1}}
 \right\}.
 \label{eq:BBC-bound}
\end{equation}

\par\medskip
\begin{theorem}[Current universal bounds]
\label{thm:known-universal-bounds}
For every \(R>1\),
\begin{align}
 2
 &\leq K(R)
 \leq
 \min\left\{
  1+\sqrt2,\,
  K_{\mathrm{Tom}}(R)
 \right\},
 \label{eq:scalar-bounds}\\
 2
 &\leq K_{\cb}(R)
 \leq
 \min\left\{
  K_{\mathrm{BBC}}(R),\,
  K_{\mathrm{Tom}}(R)
 \right\}.
 \label{eq:complete-bounds}
\end{align}
Consequently,
\begin{equation}
 \lim_{R\to\infty}K(R)
 =
 \lim_{R\to\infty}K_{\cb}(R)
 =2.
 \label{eq:asymptotic-two}
\end{equation}
\end{theorem}

\noindent
The ingredients of Theorem~\ref{thm:known-universal-bounds} are:
\begin{itemize}
 \item Tsikalas's lower bound \(K(R)\geq2\)
 \cite{Tsikalas2022};
 \item the uniform scalar estimate \(K(R)\leq1+\sqrt2\) of
 Crouzeix and Greenbaum \cite{CrouzeixGreenbaum2019};
 \item the complete intersection-of-disks estimate
 \(K_{\cb}(R)\leq K_{\mathrm{BBC}}(R)\leq2+2/\sqrt3\) of
 Badea, Beckermann, and Crouzeix \cite{BadeaBeckermannCrouzeix2009};
 \item Tomar's complete estimate
 \(K_{\cb}(R)\leq K_{\mathrm{Tom}}(R)\)
 \cite{Tomar2026}.
\end{itemize}
Pascoe first established the exact asymptotic value
\(\lim_{R\to\infty}K(R)=2\) through a degeneration to the quantum
cross \cite{Pascoe2025}.  The complete estimate
\eqref{eq:complete-bounds} now also yields the complete asymptotic
statement in \eqref{eq:asymptotic-two}.

The unresolved scalar gap can be displayed as
\begin{equation}
 0\leq K(R)-2
 \leq
 \min\left\{
  \sqrt2-1,\,
  \frac{4R^2}{(R^2+1)\sqrt{R^4-1}}
 \right\}.
 \label{eq:unresolved-gap}
\end{equation}
In particular,
\[
 K_{\mathrm{Tom}}(R)
 =
 2+\frac{4}{R^2}+O(R^{-4})
 \qquad(R\to\infty).
\]

\begin{theorem}[Finite-dimensional bounds]
\label{thm:finite-dimensional-bound}
For matrices of arbitrary finite size,
\begin{equation}
 2\leq K_{\fd}(R)
 \leq
 1+\sqrt{1+\frac{2}{R^2+1}}.
 \label{eq:Crouzeix-matrix-bound}
\end{equation}
\end{theorem}

The lower bound for the all-finite-sizes constant is recorded in
\cite{Tsikalas2026}; it does not assert that every fixed dimension has
constant at least \(2\) (indeed \(K_1(R)=1\)).  The upper bound is
Crouzeix's finite-dimensional estimate \cite{Crouzeix2025}.  It is
stated there under the strict assumptions
\(\|T\|<R\) and \(\|T^{-1}\|<R\); the non-strict version for rational
functions follows by applying the result with \(R'>R\), chosen close
enough that the function has no pole on
\(\overline{\mathbb A}_{R'}\), and then letting \(R'\downarrow R\).
The argument producing \eqref{eq:Crouzeix-matrix-bound} is not
presently known to extend to arbitrary Hilbert-space operators or to
the completely bounded setting.

\begin{remark}[Status]
The exact value at a fixed finite modulus remains open; this is
stated explicitly in the recent literature
\cite{PalPascoeTomar2026,Tsikalas2026}.  The equality
\(K(R)=2\) should therefore be described as the sharp candidate, not
as an established result.  The quantum cross has exact constant
\(2\), but it is a limiting problem and does not settle any fixed
finite \(R\) \cite{Pascoe2025}.
\end{remark}

\begin{remark}[The degenerate endpoint]
If the parameter is formally extended to \(R=1\), the operator class
consists of unitaries and the corresponding constant is \(1\).
Thus the lower bound \(K(R)\geq2\) for every \(R>1\) forces a jump at
the degenerate endpoint.  The behavior as \(R\downarrow1\) from the
nondegenerate side is itself an interesting sharp-constant question.
\end{remark}

\section{Why the problem is not two separate disk inequalities}

The two norm assumptions in \eqref{eq:quantum-annulus} separately
give contraction estimates:
\[
 \|R^{-1}T\|\leq1,
 \qquad
 \|R^{-1}T^{-1}\|\leq1.
\]
Thus the outer disk controls analytic powers of \(T\), and the inner
disk controls analytic powers of \(T^{-1}\).  A general Laurent
polynomial, however, couples these two parts.  Intersections of
spectral sets need not remain spectral sets with constant \(1\).
The quantum annulus problem asks for the optimal cost of this
coupling.

McCullough and Pascoe proved that every \(T\in\QAR\) dilates to a
quantum annulus unitary \(J\) in such a way that
\[
 T^n=P_{\mathcal H}J^n|_{\mathcal H}
 \qquad(n\in\mathbb Z)
\]
\cite{McCulloughPascoe2023}.  This is a powerful reduction, but
\(J\) need not be a normal operator supported on
\(\partial\AR\).  Hence the dilation does not by itself give the
constant \(1\), or determine the sharp constant.

Analytic proofs face the same two-boundary phenomenon.  On a convex
domain, positivity of the double-layer kernel is a central source of
spectral-set bounds.  For an annulus, the inner boundary has the
opposite orientation and produces an additional term.  A sharp proof
must preserve enough interaction between the two boundaries to reach
the lower bound \(2\), rather than estimate the two contributions
independently.

\section{Exact requirements for a complete solution}

A proof of the scalar conjecture for a fixed \(R>1\) must establish
\eqref{eq:scalar-laurent-target} uniformly in all of the following:
\[
 \text{Hilbert space dimension},\qquad
 T\in\QAR,\qquad
 \text{Laurent degree},\qquad
 \text{scalar coefficients}.
\]
Because \(K(R)\geq2\), such an upper bound would immediately imply
\(K(R)=2\).

A proof restricted to matrices would determine
\(K_{\fd}(R)\), but an additional approximation, compactness, or
dilation argument would be needed to identify that value with the
all-Hilbert-space constant.  Similarly, a scalar proof does not
automatically establish the complete constant: the latter requires
the uniform matricial inequality \eqref{eq:matrix-laurent-target}.

The following are therefore clean intermediate targets.

\begin{problem}[Finite-dimensional sharp constant]
\label{prob:finite-dimensional}
Determine \(K_n(R)\) for each \(n\), and determine
\[
 K_{\fd}(R)=\sup_{n\geq1}K_n(R).
\]
In particular, decide whether \(K_{\fd}(R)=2\).
\end{problem}

\begin{problem}[Finite-to-infinite transfer]
\label{prob:finite-infinite}
Determine whether
\[
 K_{\fd}(R)=K(R)
\]
for every \(R>1\), and, if so, construct a transfer principle that
preserves both bounds
\(\|T\|\leq R\) and \(\|T^{-1}\|\leq R\).
\end{problem}

\begin{problem}[Scalar-to-complete transfer]
\label{prob:scalar-complete}
Determine whether
\[
 K_{\cb}(R)=K(R)
\]
for every \(R>1\).  Equivalently, determine whether a sharp scalar
annulus inequality automatically amplifies to all matrix levels
without loss.
\end{problem}

\begin{problem}[Extremal structure]
\label{prob:extremals}
Classify exact or asymptotic extremizers.  In particular, determine
whether the lower value \(2\) can be approached, for fixed \(R\), by
matrices of bounded dimension, or whether increasing dimension is
intrinsic.
\end{problem}

\section{A concise master formulation}

The complete research problem may be compressed into one statement.

\begin{problem}[Master problem]
\label{prob:master}
Fix \(R>1\).  Find the least constants \(K_{\mathrm{sc}}(R)\) and
\(K_{\mathrm{mat}}(R)\) such that, for every complex Hilbert space
\(\mathcal H\), every invertible
\(T\in\mathcal B(\mathcal H)\) with
\(\|T\|,\|T^{-1}\|\leq R\), and every Laurent polynomial
\[
 F(z)=\sum_{k=-N}^{N}A_kz^k,
\]
one has
\[
 \left\|\sum_{k=-N}^{N}A_k\otimes T^k\right\|
 \leq
 K_{\mathrm{mat}}(R)
 \max_{R^{-1}\leq|z|\leq R}\|F(z)\|.
\]
Here \(A_k\in\Mat_m(\C)\), uniformly over \(m\geq1\), while the
scalar constant \(K_{\mathrm{sc}}(R)\) is obtained by restricting to
\(m=1\).  Thus
\[
 K_{\mathrm{sc}}(R)=K(R),
 \qquad
 K_{\mathrm{mat}}(R)=K_{\cb}(R).
\]
Determine both constants exactly, decide whether they coincide, and
in particular decide whether both are equal to \(2\).
\end{problem}

\begin{thebibliography}{99}

\bibitem{BadeaBeckermannCrouzeix2009}
C. Badea, B. Beckermann, and M. Crouzeix,
\emph{Intersections of several disks of the Riemann sphere as
\(K\)-spectral sets},
Commun. Pure Appl. Anal. \textbf{8} (2009), no.~1, 37--54.
\url{https://doi.org/10.3934/cpaa.2009.8.37}.

\bibitem{CrouzeixGreenbaum2019}
M. Crouzeix and A. Greenbaum,
\emph{Spectral sets: numerical range and beyond},
SIAM J. Matrix Anal. Appl. \textbf{40} (2019), no.~3,
1087--1101.
\url{https://doi.org/10.1137/18M1198417}.

\bibitem{Crouzeix2025}
M. Crouzeix,
\emph{Spectral estimates in the quantum annulus and in the numerical
annulus},
arXiv:2512.11813, 2025.
\url{https://arxiv.org/abs/2512.11813}.

\bibitem{McCulloughPascoe2023}
S. McCullough and J.~E. Pascoe,
\emph{Geometric dilations and operator annuli},
arXiv:2202.08872v4, 2023.
\url{https://arxiv.org/abs/2202.08872}.

\bibitem{PalPascoeTomar2026}
S. Pal, J.~E. Pascoe, and N. Tomar,
\emph{Spectral constants for the quantum annulus},
arXiv:2601.17560v2, 2026; Communications on Pure and Applied
Analysis, to appear.
\url{https://arxiv.org/abs/2601.17560}.

\bibitem{Pascoe2025}
J.~E. Pascoe,
\emph{The spectral constant for the quantum cross and asymptotically
sharp bounds for annuli},
arXiv:2505.06230, 2025.
\url{https://arxiv.org/abs/2505.06230}.

\bibitem{Tomar2026}
N. Tomar,
\emph{Regular quantum annulus unitary dilation and applications},
arXiv:2606.14366, 2026.
\url{https://arxiv.org/abs/2606.14366}.

\bibitem{Tsikalas2022}
G. Tsikalas,
\emph{A note on a spectral constant associated with an annulus},
Oper. Matrices \textbf{16} (2022), no.~1, 95--99.
\url{https://doi.org/10.7153/oam-2022-16-09}.

\bibitem{Tsikalas2026}
G. Tsikalas,
\emph{Spectral estimates over multiply connected domains},
arXiv:2607.10107, 2026.
\url{https://arxiv.org/abs/2607.10107}.

\end{thebibliography}

\end{document}
